\section{Model 会变好，应用公司做环境和经验}

上一章讲 Agent 生态，本章讲应用公司的边界。明超平说“永远相信 Model 会变好，永远相信 Model 和你无关”。这句话不是消极，而是给应用公司划边界：基础模型公司在造更聪明的人，应用公司要用环境和经验，让这些更聪明的模型适配真实生产需求。

\begin{figure}[H]
\centering
\includegraphics[width=0.96\textwidth]{os-agent-living-system.png}
\caption{OS Agent：它是活的。基础模型造更聪明的人，应用公司提供环境和经验。自制概念图，依据 01:49:50--01:58:26 对谈内容整理。}
\end{figure}

\begin{knowledgebox}{读图：应用公司不是模型公司的低配版本}
图中 Model 变聪明，Environment 承载使用场景，Experience 适配生产需求，OS Agent 成为活系统，Feedback 持续回流。应用公司价值在于让模型能力进入可重复的工作流和用户环境。
\end{knowledgebox}

\subsection{更高效地压榨智能}

明超平把过去两年的观察概括为更高效地消耗 token、压榨智能。Chatbot 时代，token 消耗主要依赖用户 prompt；Agent 时代，系统可以主动拆任务、调工具、调子 Agent、生成中间产物，从而让 token 消耗速度提高。问题是，这种消耗必须产生价值，否则只是成本。

\begin{importantbox}{应用公司的关键问题}
不是“我有没有模型”，而是“我能不能让模型更高效地变成用户完成任务的能力”。这要求产品理解任务、环境、反馈、分发和商业模式，而不是只把 API 包一层。
\end{importantbox}

\subsection{定价与商业模式}

前面讲环境和智能消耗，本节看它如何转成商业模式。访谈里提到 Agent 产品的定价不一定只基于服务收费。订阅、服务收费、广告、创作者激励和平台抽成都可能存在。明超平的判断是，商业模式要创造平台和创作者的共赢：公司有足够利润继续做产品，用户或创作者也能得到足够激励。

\begin{knowledgebox}{术语消化：AI 应用商业化}
\noindent\begin{tabularx}{\textwidth}{p{0.20\textwidth}p{0.34\textwidth}X}
\toprule
模式 & 适用场景 & 风险 \\
\midrule
订阅 & 高频个人或团队工具 & 容易被模型公司捆绑压价。 \\
按服务收费 & 明确任务产出，如生成网站、部署、执行流程 & 需要成本和交付质量可控。 \\
广告 & 内容和分发场景 & 可能损害创作体验。 \\
创作者分成 & 平台生态和作品交易 & 需要真实交易和版权治理。 \\
企业收费 & 团队协作、权限、安全和部署 & 销售周期长，产品会变重。 \\
\bottomrule
\end{tabularx}
\end{knowledgebox}

\subsection{为什么不是简单 to B 化}

明超平提到，如果把所有 to C 公司都变成 to B 公司，会受到大家的反抗。这句话背后是 AI 应用创业的路线问题。很多 Agent 产品为了收入确定性会向企业服务靠拢，但一旦完全 to B 化，产品可能变重、销售周期变长、创作社区变弱，也会丢掉新用户群和网络效应。

YouWare 的张力在这里：它既需要商业化，又不能过早把自己锁进传统企业软件。coding 创作、作品发布、Vibe Coder 和 Agent 网络更像 to C 或 prosumer 生态；权限、身份、团队协作和支付又会把它拉向 B 端。高标准的产品策略应该保留两边可能性，而不是为了短期收入牺牲未来入口。

\begin{importantbox}{路线选择：先成为环境，再成为销售品}
如果 YouWare 先成为创作环境，就有机会沉淀作品、用户、Agent 调用和网络效应；如果过早只做企业交付，可能获得收入，但很难形成 AgentRank 和新人群心智。
\end{importantbox}

\subsection{本章小结}

应用公司的价值不在复制基础模型，而在建立环境、经验和反馈。YouWare 的长期问题，是能否从一个 coding 创作入口，变成持续消耗智能、产生作品和形成网络效应的活系统。

